% Source: https://artofproblemsolving.com/community/c3h1883574_isosceles_legs_and_base
% Stable identity: aops-topic-1883574-q01
Two of the three sides must be equal.
\begin{itemize}
\item $x-4=2x-9$ gives $x=5$, with sides $1,1,0$. A side of length $0$ is impossible (equivalently, $1+0>1$ fails), so this is not a triangle.
\item $x-4=3x-15$ gives $x=\dfrac{11}{2}$, with sides $\dfrac{3}{2},2,\dfrac{3}{2}$. The triangle inequality holds: $\dfrac{3}{2}+\dfrac{3}{2}=3>2$.
\item $2x-9=3x-15$ gives $x=6$, with sides $2,3,3$, and $2+3>3$.
\end{itemize}
(All three sides cannot be equal, since the first two equations give different values of $x$.)

\begin{asy}
usepackage("amsmath");
size(6.5cm);
defaultpen(fontsize(10pt));
pair A=(0,0),B=(2,0),C=(1,sqrt(1.25));
draw(A--B--C--cycle,linewidth(1));
label("$2$",(A+B)/2,S);label("$\dfrac{3}{2}$",(A+C)/2+(-0.42,0.12));label("$\dfrac{3}{2}$",(B+C)/2+(0.42,0.12));
pair A2=(4,0),B2=(6,0),C2=(5,sqrt(8));
draw(A2--B2--C2--cycle,linewidth(1));
label("$2$",(A2+B2)/2,S);label("$3$",(A2+C2)/2,W);label("$3$",(B2+C2)/2,E);
label("$x=\dfrac{11}{2}$",(1,-1.0));label("$x=6$",(5,-1.0));
label("\mbox{}",(min(currentpicture,user=true)+max(currentpicture,user=true))/2,black); // keeps the text after the figure black when LaTeX typesets the labels inline
\end{asy}

So the possible values are $\dfrac{11}{2}$ and $6$, with sum $\boxed{\dfrac{23}{2}}$.
